\exercisesection{}

In the following exercises, when E is a Banach space, we denote by \(\mathcal{C}alk(E)\) the Calkin algebra \(\mathcal{L}(E)/\mathcal{L}^{c}(E)\) . We denote by \(\pi\) the canonical projection from \(\mathcal{L}(E)\) onto \(\mathcal{C}alk(E)\) .

If E is a complex Hilbert space, we equip \(\mathcal{C}alk(\mathrm{E})\) with the involution induced from the involution of \(\mathcal{L}(\mathrm{E})\) by passage to the quotient.

\begin{enumerate}
\def\labelenumi{\arabic{enumi})}
\tightlist
\item
  Let E be a locally convex space.
\end{enumerate}

\begin{enumerate}
\def\labelenumi{\alph{enumi})}
\item
  Suppose that every neighborhood of 0 in E contains a vector subspace of E of finite codimension. Show that the space \(\mathcal{F}(\mathrm{E})\) of Fredholm endomorphisms of E is not closed in \(\mathcal{L}(\mathrm{E})\) endowed with the topology of bounded convergence.
\item
  Give an example of an infinite-dimensional Fréchet space such that every neighborhood of 0 contains a vector subspace of E of finite codimension. c) Let \(E_{\sigma}\) be an infinite-dimensional Banach space equipped with the weak topology. The set \(\mathcal{F}(E_{\sigma})\) is not open in \(\mathcal{L}(E_{\sigma})\) .
\end{enumerate}

\begin{enumerate}
\def\labelenumi{\arabic{enumi})}
\setcounter{enumi}{1}
\tightlist
\item
  Let E be an infinite-dimensional Banach space and \(u \in \mathcal{L}(\mathrm{E})\) . We denote by \(\mathrm{S}(u)\) the set of complex numbers \(\lambda\) such that there exist a real number \(\delta > 0\) and a finite-codimensional subspace F of E such that \(\| u(x) - \lambda x \| \geqslant \delta \| x \|\) for all \(x \in \mathrm{F}\) .
\end{enumerate}

\begin{enumerate}
\def\labelenumi{\alph{enumi})}
\item
  The set \(S(u)\) is open in C and contains the complement in C of the spectrum of u.
\item
  If \(\lambda \in \mathrm{S}(u)\) , the endomorphism \(u - \lambda 1_{\mathrm{E}}\) has a closed image and a finite-dimensional kernel.
\item
  The complement \(\mathrm{F}(u) = \mathbf{C} - \mathrm{S}(u)\) is not empty. (Show that an element of \(\mathrm{Sp}_{\mathrm{e}}(u)\) of maximal modulus belongs to \(\mathrm{F}(u)\) .)
\end{enumerate}

\begin{enumerate}
\def\labelenumi{\arabic{enumi})}
\setcounter{enumi}{2}
\tightlist
\item
  Let E be an infinite-dimensional complex Hilbert space.
\end{enumerate}

\begin{enumerate}
\def\labelenumi{\alph{enumi})}
\item
  The involutive algebra \(\mathcal{C}alk(E)\) is a star algebra.
\item
  Let u be an element of \(\mathcal{L}(\mathrm{E})\) such that \(u^{*}u\) is compact. Show that u is compact.
\item
  Let G be the group of invertible elements in \(\mathcal{C}alk(\mathrm{E})\) and \(G_{0}\) its neutral component. Let \(p: G \to G/G_{0}\) be the canonical projection. There exists a unique group isomorphism \(\alpha: Z \to G/G_{0}\) such that \(\alpha(\operatorname{ind}(u)) = p(\pi(u))\) for every Fredholm endomorphism u of E.
\end{enumerate}

\begin{enumerate}
\def\labelenumi{\arabic{enumi})}
\setcounter{enumi}{3}
\tightlist
\item
  Let E be a complex Hilbert space. For every invertible element \(\xi\) of \(\mathcal{C}alk(E)\) , we call the index of \(\xi\) the index of any Fredholm endomorphism \(u\) of E such that \(\pi(u) = \xi\) .
\end{enumerate}

\begin{enumerate}
\def\labelenumi{\alph{enumi})}
\item
  Show that for every \(\xi \in \mathcal{C}alk(\mathrm{E})\) that is Hermitian, there exists \(u \in \mathcal{L}(\mathrm{E})\) that is Hermitian such that \(\pi(u) = \xi\) .
\item
  Show that for every \(\xi \in \mathcal{C}alk(\mathrm{E})\) that is unitary and has index zero, there exists \(u \in \mathcal{L}(\mathrm{E})\) that is unitary such that \(\pi(u) = \xi\) . (Let \(v\) be invertible such that \(\pi(v) = \xi\) ; consider the polar decomposition \(v = j|v|\) of \(v\) and verify that \(|v| - 1_{\mathrm{E}}\) is compact, then deduce that \(v - j\) is compact.)
\item
  Show that for every \(\xi \in \mathcal{C}alk(\mathrm{E})\) that is unitary such that \(\xi^2 = 1\) , there exists \(u \in \mathcal{L}(\mathrm{E})\) that is unitary such that \(u^2 = 1\) and \(\pi(u) = \xi\) .
\end{enumerate}

\(d)\) For every element \(p \in \mathcal{C}alk(\mathrm{E})\) such that \(p^2 = p\) , there exists a projector \(\mathrm{P} \in \mathcal{L}(\mathrm{E})\) such that \(\pi(\mathrm{P}) = p\) .

\begin{enumerate}
\def\labelenumi{\arabic{enumi})}
\setcounter{enumi}{4}
\tightlist
\item
  Let E be a complex Hilbert space.
\end{enumerate}

\begin{enumerate}
\def\labelenumi{\alph{enumi})}
\tightlist
\item
  Let \(n\) be an integer with \(n \geqslant 2\) . We set \(\mathrm{I} = \mathbf{N} \times \{0, \ldots, n\}\) . Let \((\sigma_i)_{1 \leqslant i \leqslant n}\) be the family of permutations of I defined by
\end{enumerate}

\[
\begin{array}{c} \sigma_ {i} (m, i) = (m - 1, i) \text {for} m \geqslant 1 \\ \sigma_ {i} (0, i) = (0, i - 1) \\ \sigma_ {i} (m, i - 1) = (m + 1, i - 1) \text {for} m \geqslant 0 \\ \sigma_ {i} (m, j) = (m, j) \text {for} m \geqslant 0 \text {and} j \notin \{i, i - 1 \}. \end{array}
\]

\begin{enumerate}
\def\labelenumi{\alph{enumi})}
\setcounter{enumi}{1}
\item
  We denote by \((e_{m,i})\) the canonical orthonormal basis of \(\ell^{2}(\mathrm{I})\) . Let \(u_{i}\) be the endomorphism of \(\ell^{2}(\mathrm{I})\) such that \(u_{i}(e_{m,j}) = e_{\sigma_{i}(m,j)}\) for all \((m,j) \in \mathrm{I}\) . The endomorphism \(u_{i}\) is unitary.
\item
  The elements \(\pi(u_i)\) are pairwise commuting.
\end{enumerate}

\(d)\) There does not exist a family \((v_{1},\ldots,v_{n})\) of unitary endomorphisms of E, pairwise commuting, such that \(\pi(v_{i})=\pi(u_{i})\) for all \(i\) .

\begin{enumerate}
\def\labelenumi{\arabic{enumi})}
\setcounter{enumi}{5}
\tightlist
\item
  Let E be a complex Hilbert space. Recall that \(\mathcal{C}alk(E)\) is a star algebra (exercise 3).
\end{enumerate}

  Let \(\mathfrak{F}\) be a nonprincipal ultrafilter on \(\mathbf{N}\) . Let F be the vector space of bounded sequences in E that converge weakly to 0 along \(\mathfrak{F}\) and let \(F_0\) be the subspace of F consisting of sequences that converge in norm to 0 along \(\mathfrak{F}\) . We set \(\widetilde{\mathrm{H}} = \mathrm{F} / \mathrm{F}_0\) . a) For all elements \(x = (x_n)\) and \(y = (y_n)\) of F, show that the sequence \((\langle x_n | y_n \rangle)_{n \in \mathbf{N}}\) converges along \(\mathfrak{F}\) . We denote by \(\langle x | y \rangle_{\mathrm{F}}\) its limit.

\begin{enumerate}
\def\labelenumi{\alph{enumi})}
\setcounter{enumi}{1}
\item
  Show that \((x,y)\mapsto\langle x\mid y\rangle_{\mathrm{F}}\) is a positive Hermitian form on F, and that \(\langle x\mid x\rangle_{F}=0\) if and only if \(x\in F_{0}\) ; this Hermitian form therefore defines a separated pre-Hilbert space structure on \(\widetilde{H}\).
\item
  Let H be the Hilbert space obtained by completing the separated pre-Hilbert space \(\widetilde{H}\) . Show that for every u in \(\mathcal{L}(E)\) , one obtains an element \(\varrho(u)\) of \(\mathcal{L}(H)\) by passing to the quotient of the linear map \((x_{n})\mapsto(u(x_{n}))\) from F to F. d) The map \(u\mapsto\varrho(u)\) induces an isometric morphism of C*-algebras \(\mathcal{C}alk(E)\to\mathcal{L}(H)\) .
\end{enumerate}

(We will see later that every C*-algebra A has an isometric representation \(A \to \mathcal{L}(E)\) for a certain Hilbert space E, cf. th. 2 of V, p. 442.)

\begin{enumerate}
\def\labelenumi{\arabic{enumi})}
\setcounter{enumi}{6}
\tightlist
\item
  We equip the unit circle \(S_{1} \subset \mathbf{C}\) with the normalized Haar measure. Let E be the complex Hilbert space \(\mathrm{L}^{2}(\mathbf{S}_{1})\) . For \(n \in \mathbf{Z}\) , let \(e_{n}\) be the element \(z \mapsto z^{n}\) of E. Let F ⊂ E be the closed subspace generated by the \(e_{n}\) for \(n \geqslant 0\) and let \(p: E \to F\) be the orthogonal projector. For \(f \in L^{\infty}(\mathbf{S}_{1})\) we define linear maps \(M_{f}\) (resp. \(T_{f}\) ) on E (resp. F) by \(\mathrm{M}_{f}(g) = fg\) for \(g \in E\) and \(T_{f} = p \circ M_{f}\) . One says that \(T_{f}\) is the Toeplitz operator with symbol f. (a) The linear maps \(M_{f}\) and \(T_{f}\) are continuous. The spectrum of \(M_{f}\) is the closure of the set of Fourier coefficients \(c_{n} = \langle f | e_{n} \rangle\) of f. The endomorphism \(M_{f}\) is a Fredholm endomorphism if and only if it is a Riesz endomorphism.
\end{enumerate}

\begin{enumerate}
\def\labelenumi{\alph{enumi})}
\setcounter{enumi}{1}
\item
  The endomorphisms \(T_{e_{n}}\) for \(n \in Z\) are Fredholm endomorphisms and \(\text{ind}(\mathrm{T}_{e_{n}}) = -n\) for all \(n \in Z\) .
\item
  One has \(\mathrm{T}_{\overline{f}} = \mathrm{T}_f^*\) for every function \(f\in \mathrm{L}^{\infty}(\mathbf{S}_1)\) .
\item
  If \(T_{f}\) is an isomorphism, then \(M_{f}\) is an isomorphism and f is invertible in \(\mathrm{L}^{\infty}(\mathbf{S}_{1})\) . Deduce from this that the spectrum of \(M_{f}\) is contained in the spectrum of \(T_{f}\) and that \(\|T_{f}\|=\|f\|\) .
\item
  Let \(f \in \mathcal{C}(\mathbf{S}_1)\) be a continuous function and \(g \in \mathrm{L}^{\infty}(\mathbf{S}_1)\) . The endomorphisms \(\mathrm{T}_f\mathrm{T}_g - \mathrm{T}_{fg}\) and \(\mathrm{T}_g\mathrm{T}_f - \mathrm{T}_{gf}\) are compact. (First consider the case where \(f = e_n\) .)
\item
  Let \(\mathcal{C}alk(\mathrm{F})\) be the Calkin algebra of F and \(\pi\colon\mathcal{L}(\mathrm{F})\to\mathcal{C}alk(\mathrm{F})\) be the canonical projection. Show that \(f\mapsto\pi(\mathrm{T}_{f})\) is an injective morphism of involutive algebras from \(\mathcal{C}(\mathbf{S}_{1})\) into \(\mathcal{C}alk(\mathrm{F})\) (To prove injectivity, one may use prop. 1 of I, p. 31.)
\end{enumerate}

\(g)\) Let \(f \in \mathcal{C}(\mathbf{S}_1)\) . Then \(\mathrm{T}_f\) is a Fredholm endomorphism if and only if \(f\) does not vanish on \(\mathbf{S}_1\) . (Use prop. 5 of I, p. 106.)

\(h)\) Let \(f \in \mathcal{C}(\mathbf{S}_1)\) such that \(f\) does not vanish on \(\mathbf{S}_1\) . The index of \(\mathrm{T}_f\) is equal to \(n \in \mathbf{Z}\) if and only if the path \(t \mapsto f(e^{2i\pi t})\) in \(\mathbf{C} - \{0\}\) is homotopic to the path \(t \mapsto e_n(t)\) (cf. TA, III, p. 229, def. 1).

\begin{enumerate}
\def\labelenumi{\arabic{enumi})}
\setcounter{enumi}{7}
\tightlist
\item
  We resume the notation of the preceding exercise.
\end{enumerate}

\begin{enumerate}
\def\labelenumi{\alph{enumi})}
\item
  We denote by \(\mathrm{T}(\mathbf{S}_{1})\) the closed subalgebra of \(\mathcal{L}(\mathrm{F})\) generated by the Toeplitz operators \(T_{e_{1}}\) and \(T_{e_{-1}}\) . Show that \(\mathrm{T}(\mathbf{S}_{1})\) is a star subalgebra of \(\mathcal{L}(\mathrm{F})\) .
\item
  Let \(f \in \mathcal{C}(\mathbf{S}_1)\) . Show that \(T_f\) belongs to \(T(\mathbf{S}_1)\) .
\item
  The representation of \(\mathrm{T}(\mathbf{S}_{1})\) on F is irreducible; the star algebra \(\mathrm{T}(\mathbf{S}_{1})\) contains \(\mathcal{L}^{\mathrm{c}}(\mathrm{F})\) .
\end{enumerate}

\(d)\) One has \(u \in \mathrm{T}(\mathbf{S}_1)\) if and only if there exists \(f \in \mathcal{C}(\mathbf{S}_1)\) and \(h \in \mathcal{L}^c(\mathrm{F})\) such that \(u = \mathrm{T}_f + h\) . (Show that the set of endomorphisms of the form \(\mathrm{T}_f + h\) forms a closed star algebra in \(\mathcal{L}(\mathrm{F})\) .)

\begin{enumerate}
\def\labelenumi{\alph{enumi})}
\setcounter{enumi}{4}
\tightlist
\item
  There exists an exact sequence of morphisms of involutive algebras
\end{enumerate}

\[
0 \longrightarrow \mathcal {L} ^ {\mathrm{c}} (\mathrm{F}) \longrightarrow \mathrm{T} (\mathbf {S} _ {1}) \stackrel {\varrho} {\longrightarrow} \mathcal {C} (\mathbf {S} _ {1}) \longrightarrow 0
\]

such that \(\varrho(\mathrm{T}_f) = f\) for every function \(f \in \mathcal{C}(\mathbf{S}_1)\) .

\begin{enumerate}
\def\labelenumi{\arabic{enumi})}
\setcounter{enumi}{8}
\item
  Let B be an infinite-dimensional Banach space. We denote by E the strong dual of B and by F the dual of B equipped with the topology \(\sigma(\mathrm{B}^{\prime},\mathrm{B})\) . Show that there exists a continuous bijective linear map u from E into F and a compact linear map h from E into F such that \(u + h\) has a cokernel of infinite dimension.
\item
  Let E be a separated locally convex space, and let u and h be commuting elements of \(\mathcal{L}(\mathrm{E})\) . Suppose there exists an integer n such that \(h^{n}\) is a compact endomorphism.
\end{enumerate}

\begin{enumerate}
\def\labelenumi{\alph{enumi})}
\item
  If u is a Fredholm endomorphism, then u-h is a Fredholm endomorphism.
\item
  If \(u\) is a Riesz endomorphism, then \(u - h\) is a Riesz endomorphism.
\end{enumerate}

\begin{enumerate}
\def\labelenumi{\arabic{enumi})}
\setcounter{enumi}{10}
\tightlist
\item
  Let \(\mathrm{E} = \ell^2 (\mathbf{N})\) and let \((e_n)\) be the canonical basis of E. We define \(u\in \mathcal{L}(\mathrm{E})\) by \(u(e_i) = e_{i + 1}\) for \(i\in \mathbf{N}\) . Let \(\mathrm{F} = \mathrm{E}\oplus \mathrm{E}\) and \(v = u\oplus (u^{*} + 2\cdot 1_{\mathrm{E}})\in \mathcal{L}(\mathrm{F})\) . Let \(\overline{v} = \pi (v)\) be the image of \(v\) in the Calkin algebra \(\mathcal{Calk}(\mathrm{F})\) .
\end{enumerate}

\begin{enumerate}
\def\labelenumi{\alph{enumi})}
\item
  The endomorphism \(v\) is a Fredholm endomorphism of index -1, and \(v - 2 \cdot 1_{\mathrm{E}}\) is a Fredholm endomorphism of index 1.
\item
  Let \(p = \mathrm{X}(\mathrm{X} - 2) \in \mathbf{C}[\mathrm{X}]\) and let S be the exponential spectrum of \(\overline{v}\) (exercise 5 of I, p. 173). We have \(0 \in p(\mathrm{S})\) but 0 does not belong to the exponential spectrum of \(p(\overline{v})\) (Use Exercise 3.)
\end{enumerate}

